\begin{drillsolution}{d:34}\textbf{逐段追踪入射高度和倾角}\par\nobreak
图\ref{fig:d34q}中取第一镜处入射高度为 $y_0>0$、入射倾角为零。经过第一镜，倾角变为 $\alpha_1=-y_0/f_1$；传播到第二镜时，高度为 $y_2=y_0+d\alpha_1=y_0(1-d/f_1)$。第二镜再改变倾角，故
\[
\alpha_2=\alpha_1-\frac{y_2}{f_2}=-\Phi_{\rm eq}y_0,
\qquad \Phi_{\rm eq}=\frac1{f_1}+\frac1{f_2}-\frac d{f_1f_2}.
\]
第二镜后的轴交点满足 $0=y_2+b\alpha_2$，其中 $b$ 向右为正。因此在 $\Phi_{\rm eq}\ne0$ 时，
\[
b=\frac{1-d/f_1}{\Phi_{\rm eq}},\qquad f_{\rm eq}=\frac1{\Phi_{\rm eq}}.
\]

\par\medskip\Needspace{4\baselineskip}\noindent\textbf{分别判断两种镜距下的实际传播}\par\nobreak
$d=15.0\,\mathrm{cm}$ 时，$\Phi_{\rm eq}=0.0750\,\mathrm{cm}^{-1}$，$f_{\rm eq}=13.33\,\mathrm{cm}$，但 $y_2=-y_0/2$、$\alpha_2<0$。图\ref{fig:d34s}上图的光线已在两镜之间越过轴线，出第二镜后继续离轴，故为发散束，其反向延长线交于虚后焦点：
\[
\boxed{b=-6.667\,\mathrm{cm}}.
\]
$d=30.0\,\mathrm{cm}$ 时，$\Phi_{\rm eq}=0$、$y_2=-2y_0$、$\alpha_2=0$，图中下方出射光与轴平行，系统无焦，没有有限后焦点。
\fig{fig:d34s}{上图在第二镜前已越轴，出射束发散；下图重新成为平行束}{\begin{tikzpicture}
\begin{scope}
\draw[axisline](-1,0)--(6.3,0);\draw[world](0,-1.8)--(0,1.6)(3,-1.8)--(3,1.6);
\draw[ray](-1,1)--(0,1);\draw[ray](0,1)--(3,-.5);\draw[ray](3,-.5)--(6,-1.625);
\draw[virtual](1.666667,0)--(3,-.5);\fill(1.666667,0)circle(1.3pt)node[above left]{$F'$};
\node[above]at(0,1.6){$L_1$};\node[above]at(3,1.6){$L_2$};\node at(5,1){$d=15.0\,\mathrm{cm}$};
\draw[dim](1.666667,-2.1)--(3,-2.1)node[midway,below]{$|b|$};
\end{scope}
\begin{scope}[yshift=-5cm]
\draw[axisline](-1,0)--(7.6,0);\draw[world](0,-2.4)--(0,1.5)(6,-2.4)--(6,1.5);
\draw[ray](-1,1)--(0,1);\draw[ray](0,1)--(6,-2);\draw[ray](6,-2)--(7.5,-2);
\node[above]at(0,1.5){$L_1$};\node[above]at(6,1.5){$L_2$};\node at(3.7,.8){$d=30.0\,\mathrm{cm}$};\node[right]at(6.5,-1.3){$\alpha_2=0$};
\end{scope}
\end{tikzpicture}}

\par\medskip\Needspace{4\baselineskip}\noindent\textbf{区分主平面焦距与顶点后焦距}\par\nobreak
等效焦距从相应主平面量起，$b$ 从第二镜平面量起。分离光组的主平面通常不在第二镜处，二者不能直接等同。图\ref{fig:d34s}上图即使 $f_{\rm eq}>0$，实际出射束仍可发散；判断实虚应结合出射高度和倾角，不能仅看等效焦距的符号。
\end{drillsolution}

\begin{drillsolution}{d:35}\textbf{分别计入两个折射面的约化光焦度}\par\nobreak
图\ref{fig:d35q}中的两个外部折射率不同。薄透镜的两面可在同一顶点平面相加，约化光焦度为
\[
\Phi=\frac{n_\ell-n_1}{R_{1s}}+\frac{n_2-n_\ell}{R_{2s}}
=\frac{0.50}{10}+\frac{4/3-1.50}{-10}
=\frac1{15}\,\mathrm{cm}^{-1}.
\]
近轴成像关系是 $n_1/u+n_2/v=\Phi$，两边不同的折射率必须保留。

\par\medskip\Needspace{4\baselineskip}\noindent\textbf{用平行入射与平行出射分别确定两个焦距}\par\nobreak
令 $u\to\infty$，得到像方焦距 $f'=n_2/\Phi=20.0\,\mathrm{cm}$；令 $v\to\infty$，得到物方焦距大小 $f=n_1/\Phi=15.0\,\mathrm{cm}$。所以物方焦点在左侧 $15.0\,\mathrm{cm}$，像方焦点在右侧 $20.0\,\mathrm{cm}$。

\par\medskip\Needspace{4\baselineskip}\noindent\textbf{计算有限物距的像并用光线核验}\par\nobreak
代入 $u=30.0\,\mathrm{cm}$：
\[
\frac1{30}+\frac{4/3}{v}=\frac1{15},\qquad
\boxed{v=40.0\,\mathrm{cm}},\qquad
\boxed{\beta=-\frac{n_1v}{n_2u}=-1}.
\]
\fig{fig:d35s}{平行入射线经像方焦点；过顶点的近轴线也因外部介质不同而折转}{\begin{tikzpicture}[y=2cm]
\fill[cyan!7](0,-.65)rectangle(4.6,.9);\draw[axisline](-3.6,0)--(4.7,0);\draw[world](0,-.65)--(0,.9);
\draw[axis,thick](-3,0)--(-3,.4);\draw[axis,thick](4,0)--(4,-.4);
\draw[ray](-3,.4)--(0,.4);\draw[ray](0,.4)--(4,-.4);\draw[ray](-3,.4)--(0,0);\draw[ray](0,0)--(4,-.4);
\node[above]at(-3,.4){$B$};\node[below]at(-3,0){$A$};\node[above]at(4,0){$A'$};\node[below]at(4,-.4){$B'$};
\fill(2,0)circle(1.2pt)node[above]{$F'$};\fill(-1.5,0)circle(1.2pt)node[below]{$F$};\node[above right]at(0,.9){$O$};
\end{tikzpicture}}
图\ref{fig:d35s}中两条实际光线在右侧相交，像倒立且等大。经过顶点的光线也不能无条件画成直线：薄镜顶点处的净关系为 $n_1\alpha_{\mathrm{in}}=n_2\alpha_{\mathrm{out}}$，所以出射倾角是入射倾角的 $3/4$。这与两个焦距不相等的结果一致。
\end{drillsolution}
